\section{Architecture}\label{sec:architecture}

\subsection{Two-Party Protocol Execution Model}

Every two-party test binary accepts three command-line arguments:
\begin{lstlisting}
./test_binary r=<role> p=<port> [ip=<address>]
\end{lstlisting}
where \texttt{r=1} designates ALICE (server) and \texttt{r=2} designates BOB (client).
The server opens a TCP socket via NetIO (\texttt{Utils/net\_io\_channel.h}), calls \texttt{listen()}, and performs exactly one \texttt{accept()}.
After accepting the client connection, the listening socket is closed --- this is a critical design constraint that affects how scripts detect server readiness (Section~\ref{sec:portdetect}).

The protocol proceeds as:
\begin{enumerate}[nosep]
  \item BOB generates all SEAL keys: secret key, public key, relinearization keys, and Galois keys (rotation keys for BSGS schedule).
  \item BOB sends public material to ALICE over the NetIO TCP connection using \texttt{FILE*}-buffered I/O with \texttt{\_IOFBF} (full buffering).
  \item Both parties execute the HE/MPC computation. For tests involving OT (e.g., \texttt{test\_fixpoint}), an \texttt{IKNPOTPack} with 4 threads is initialized.
  \item Results are reconstructed on BOB's side.
\end{enumerate}

\subsection{Single-Container Mode}

For benchmarking on a single machine, both ALICE and BOB run within the same Docker container on \texttt{localhost}:
\begin{lstlisting}
# Server in background
./test_binary r=1 p=1234 &
# Client connects
./test_binary r=2 p=1234 ip=127.0.0.1
\end{lstlisting}
This eliminates network variability and produces stable, low-variance measurements suitable for comparing relative performance across tests.

\subsection{Port Detection}\label{sec:portdetect}

A subtle but critical issue: the standard approach to checking whether a TCP port is open (\texttt{nc -z host port}) works by opening a TCP connection and immediately closing it.
Because NetIO performs exactly one \texttt{accept()}, this ``probe'' connection consumes the server's accept slot, causing the real client to receive a connection-refused error.

Our scripts detect server readiness by reading \texttt{/proc/net/tcp} directly:
\begin{lstlisting}[language=bash]
hex_port=$(printf '%04X' "$PORT")
grep -q ":${hex_port} 00000000:0000 0A" /proc/net/tcp
\end{lstlisting}
The \texttt{0A} state code indicates TCP LISTEN.
This is a read-only operation on the kernel's socket table and does not consume the accept.

\subsection{Docker Image Structure}

The CPU image (\texttt{blb:v1.0}) is built from \texttt{ubuntu:22.04} and performs a multi-stage dependency build:

\begin{table}[h]
\centering
\caption{Docker image build stages}
\label{tab:docker-stages}
\begin{tabular}{llr}
\toprule
Stage & Component & Approx.\ Time \\
\midrule
1 & System packages (build-essential, cmake, libssl, etc.) & 30\,s \\
2 & Apply patches (SEAL, emp-tool, emp-ot) & 5\,s \\
3 & Build Intel HEXL 1.2.6 & 2\,min \\
4 & Build Microsoft SEAL 4.1.2 (with HEXL) & 3\,min \\
5 & Build BLB (all libraries + test binaries) & 2\,min \\
\bottomrule
\end{tabular}
\end{table}

Scripts are copied from \texttt{docker/v1.0/scripts/} to \texttt{/docker-scripts/} inside the image.
The \texttt{ENTRYPOINT} is set to \texttt{entrypoint.sh}, which dispatches on the Docker CMD argument (\texttt{test}, \texttt{benchmark}, \texttt{manual}, \texttt{build}, \texttt{bash}, etc.).

A GPU variant (\texttt{blb:v1.0-gpu}) extends this with CUDA 12.2 and PhantomFHE, enabling three additional test binaries (\texttt{test\_HE}, \texttt{testHEGPU}, \texttt{testFFN}).
This image was prepared but could not be validated due to the absence of an NVIDIA GPU on the testbed.
